\section{Foundations}\label{sec:foundations}

\subsection{Numbers and evaluation}
In Lean~4~\cite{deMouraUllrich2021}, numerals are typed terms: $\lgN$
(\texttt{Nat}), $\lgZ$ (\texttt{Int}), $\lgQ$ (\texttt{Rat}), and so on.
Literals are interpreted by the \texttt{OfNat} typeclass.

\begin{lstlisting}
#eval 2 + 3 * 4            -- 14
#eval (2 : Rat) / 3        -- 2/3
#check Nat.add_comm        -- forall n m : Nat, n + m = m + n
\end{lstlisting}

\subsection{Definitions: functions, inductive types, structures}
Structural recursion is checked automatically by the termination checker.
For example, factorial with $0! := 1$ and $(n+1)! := (n+1)\,n!$:

\begin{lstlisting}
def factorial : Nat -> Nat
  | 0 => 1
  | n + 1 => (n + 1) * factorial n

#eval factorial 5          -- 120
\end{lstlisting}

\texttt{inductive} introduces inductive types and \texttt{structure}
introduces record types; both are represented as inductive constructions in
the CIC.

\subsection{Typeclasses}
Typeclasses dispatch interface implementations by type. For instance,
$\lgF{p}$ carries a \texttt{Field} structure whenever $p$ is prime:

\begin{lstlisting}
instance : Fact (Nat.Prime 7) := Fact.mk (by decide)
#synth Add Nat
#synth Field (ZMod 7)

structure WrappedNat where n : Nat
instance : Add WrappedNat where
  add a b := WrappedNat.mk (a.n + b.n)
\end{lstlisting}

\subsection{Exploration commands}
\texttt{\#check} (inspect a type), \texttt{\#eval} (compile and evaluate),
\texttt{\#reduce} (reduce to weak-head normal form), \texttt{\#print}
(inspect a definition), and \texttt{\#synth} (search a typeclass instance)
are the five most frequently used commands in the Lean workflow.
